% GENERATED FILE. Edit canonical lessons, then run npm run generate:books.
% canonical-source-hash: fnv1a64:cbcf2e56949c9f41
% canonical-lessons: UR-C410-imarat, UR-C410-maqam, UR-C410-bord, UR-C410-khamba, UR-C410-dikhana

\chapter{Landmarks}
\label{ch:ur-landmarks}
\hlchaptermodality{410}

\begin{chapteropening}
\textbf{By the end of this chapter:} \emph{I can give directions by landmarks: a building, a sign, a pole.}

Complete the last lesson of chapter 410: I can give directions by landmarks: a building, a sign, a pole.
\end{chapteropening}

\section[\texorpdfstring{\ur{عمارت} (imārat)}{عمارت (imārat)}]{\ur{عمارت} (imārat) — a building}

\label{lesson:UR-C410-imarat}



\begin{quote}
Before the new one: say the Urdu for a nook, then the Urdu for to peep.
\end{quote}

\subsection*{You'll want to know: \ur{عمارت}}
\textbf{\ur{عمارت}} — \emph{imārat} — \textquotedblleft{}a building\textquotedblright{}.

\begin{quote}
\ur{بینک} \ur{ایک} \ur{اونچی} \ur{عمارت} \ur{میں} \ur{ہے۔}
\end{quote}

\begin{quote}
\emph{baink ek ū̃chī imārat meṅ hai.}
\end{quote}

\begin{quote}
The bank is in a tall building.
\end{quote}

Feminine.

\begin{cousinweb}
\textbf{\ur{عمارت}} (\emph{imārat}) is usually traced, through Persian, to Arabic \emph{ʿimāra}, \textquotedblleft{}a building, a thing built\textquotedblright{}.
\end{cousinweb}

\begin{grammarlens}[title={What you've built}]
One more word for getting around.
\end{grammarlens}

\subsection*{Your turn}
\begin{itemize}
\raggedright
  \item \emph{Say it:} \emph{imārat}
  \item \emph{Say it:} \emph{imārat}, once more
  \item \emph{Say it:} \emph{baink ek ū̃chī imārat meṅ hai.} — The bank is in a tall building.
  \item \emph{Recall:} say the Urdu for a nook, then the Urdu for to peep, then say \emph{imārat} again
\end{itemize}

\begin{tcolorbox}[breakable,colback=teal!4,colframe=teal!35!black,title={Before you move on}]
What does \ur{عمارت} mean? (\textquotedblleft{}A building\textquotedblright{}.) Say it once more.
\end{tcolorbox}
\section[\texorpdfstring{\ur{مقام} (maqām)}{مقام (maqām)}]{\ur{مقام} (maqām) — a spot; a place}

\label{lesson:UR-C410-maqam}



\begin{quote}
Before the new one: say the Urdu for to peep, then the Urdu for a building.
\end{quote}

\subsection*{You'll want to know: \ur{مقام}}
\textbf{\ur{مقام}} — \emph{maqām} — \textquotedblleft{}a spot; a place\textquotedblright{}.

\begin{quote}
\ur{یہ} \ur{مقام} \ur{پہاڑوں} \ur{کے} \ur{بیچ} \ur{ہے۔}
\end{quote}

\begin{quote}
\emph{yih maqām pahāṛoṅ ke bīch hai.}
\end{quote}

\begin{quote}
This spot lies among the mountains.
\end{quote}

Masculine.

\begin{cousinweb}
\textbf{\ur{مقام}} (\emph{maqām}) is usually traced, through Persian, to Arabic \emph{maqām}, \textquotedblleft{}a place where one stands\textquotedblright{}; \textbf{\ur{مقامی}}, \textquotedblleft{}local\textquotedblright{}, is built on it.
\end{cousinweb}

\begin{grammarlens}[title={What you've built}]
One more word for getting around.
\end{grammarlens}

\subsection*{Your turn}
\begin{itemize}
\raggedright
  \item \emph{Say it:} \emph{maqām}
  \item \emph{Say it:} \emph{maqām}, once more
  \item \emph{Say it:} \emph{yih maqām pahāṛoṅ ke bīch hai.} — This spot lies among the mountains.
  \item \emph{Recall:} say the Urdu for to peep, then the Urdu for a building, then say \emph{maqām} again
\end{itemize}

\begin{tcolorbox}[breakable,colback=teal!4,colframe=teal!35!black,title={Before you move on}]
What does \ur{مقام} mean? (\textquotedblleft{}A spot; a place\textquotedblright{}.) Say it once more.
\end{tcolorbox}
\section[\texorpdfstring{\ur{بورڈ} (borḍ)}{بورڈ (borḍ)}]{\ur{بورڈ} (borḍ) — a signboard; a board}

\label{lesson:UR-C410-bord}



\begin{quote}
Before the new one: say the Urdu for a building, then the Urdu for a spot.
\end{quote}

\subsection*{You'll want to know: \ur{بورڈ}}
\textbf{\ur{بورڈ}} — \emph{borḍ} — \textquotedblleft{}a signboard; a board\textquotedblright{}.

\begin{quote}
\ur{سڑک} \ur{کے} \ur{کنارے} \ur{ایک} \ur{بڑا} \ur{بورڈ} \ur{ہے۔}
\end{quote}

\begin{quote}
\emph{saṛak ke kināre ek baṛā borḍ hai.}
\end{quote}

\begin{quote}
There's a big signboard by the road.
\end{quote}

Masculine.

\begin{cousinweb}
\textbf{\ur{بورڈ}} (\emph{borḍ}) is borrowed from English \emph{board}.
\end{cousinweb}

\begin{grammarlens}[title={What you've built}]
One more word for getting around.
\end{grammarlens}

\subsection*{Your turn}
\begin{itemize}
\raggedright
  \item \emph{Say it:} \emph{borḍ}
  \item \emph{Say it:} \emph{borḍ}, once more
  \item \emph{Say it:} \emph{saṛak ke kināre ek baṛā borḍ hai.} — There's a big signboard by the road.
  \item \emph{Recall:} say the Urdu for a building, then the Urdu for a spot, then say \emph{borḍ} again
\end{itemize}

\begin{tcolorbox}[breakable,colback=teal!4,colframe=teal!35!black,title={Before you move on}]
What does \ur{بورڈ} mean? (\textquotedblleft{}A signboard; a board\textquotedblright{}.) Say it once more.
\end{tcolorbox}
\section[\texorpdfstring{\ur{کھمبا} (khambā)}{کھمبا (khambā)}]{\ur{کھمبا} (khambā) — a pole; a pillar}

\label{lesson:UR-C410-khamba}



\begin{quote}
Before the new one: say the Urdu for a spot, then the Urdu for a signboard.
\end{quote}

\subsection*{You'll want to know: \ur{کھمبا}}
\textbf{\ur{کھمبا}} — \emph{khambā} — \textquotedblleft{}a pole; a pillar\textquotedblright{}.

\begin{quote}
\ur{بجلی} \ur{کے} \ur{کھمبے} \ur{کے} \ur{پاس} \ur{بائیں} \ur{مڑیے۔}
\end{quote}

\begin{quote}
\emph{bijlī ke khambe ke pās bāeṅ muṛiye.}
\end{quote}

\begin{quote}
Turn left by the electricity pole.
\end{quote}

Masculine.

\begin{cousinweb}
\textbf{\ur{کھمبا}} (\emph{khambā}) is inherited: CDIAL lists the Hindi cousin \emph{khambh}, \textquotedblleft{}a pillar, a pole\textquotedblright{}, under Sanskrit \emph{skambha}, \textquotedblleft{}a prop\textquotedblright{} (CDIAL 13639).
\end{cousinweb}

\begin{grammarlens}[title={What you've built}]
One more word for getting around.
\end{grammarlens}

\subsection*{Your turn}
\begin{itemize}
\raggedright
  \item \emph{Say it:} \emph{khambā}
  \item \emph{Say it:} \emph{khambā}, once more
  \item \emph{Say it:} \emph{bijlī ke khambe ke pās bāeṅ muṛiye.} — Turn left by the electricity pole.
  \item \emph{Recall:} say the Urdu for a spot, then the Urdu for a signboard, then say \emph{khambā} again
\end{itemize}

\begin{tcolorbox}[breakable,colback=teal!4,colframe=teal!35!black,title={Before you move on}]
What does \ur{کھمبا} mean? (\textquotedblleft{}A pole; a pillar\textquotedblright{}.) Say it once more.
\end{tcolorbox}
\section[\texorpdfstring{\ur{دکھانا} (dikhānā)}{دکھانا (dikhānā)}]{\ur{دکھانا} (dikhānā) — to show}

\label{lesson:UR-C410-dikhana}



\begin{quote}
Before the new one: say the Urdu for a signboard, then the Urdu for a pole.
\end{quote}

\subsection*{You'll want to know: \ur{دکھانا}}
\textbf{\ur{دکھانا}} — \emph{dikhānā} — \textquotedblleft{}to show\textquotedblright{}.

\begin{quote}
\ur{ایک} \ur{راہگیر} \ur{نے} \ur{ہمیں} \ur{راستہ} \ur{دکھایا۔}
\end{quote}

\begin{quote}
\emph{ek rāhgīr ne hameṅ rāstā dikhāyā.}
\end{quote}

\begin{quote}
A passer-by showed us the way.
\end{quote}

Past, with \ur{نے} (he or she …): \ur{اس} \ur{نے} \ur{نقشہ} \ur{دکھایا}, \emph{us ne naqshā dikhāyā}.

\begin{cousinweb}
\textbf{\ur{دکھانا}} (\emph{dikhānā}) is built on \textbf{\ur{دیکھنا}}, \textquotedblleft{}to see\textquotedblright{}, which comes from \emph{dekṣati}, a form CDIAL reconstructs (CDIAL 6507): it is the form for making someone see.
\end{cousinweb}

\begin{grammarlens}[title={What you've built}]
One more word for getting around.
\end{grammarlens}

\subsection*{Your turn}
\begin{itemize}
\raggedright
  \item \emph{Say it:} \emph{dikhānā}
  \item \emph{Say it:} \emph{dikhānā}, once more
  \item \emph{Say it:} \emph{ek rāhgīr ne hameṅ rāstā dikhāyā.} — A passer-by showed us the way.
  \item \emph{Recall:} say the Urdu for a signboard, then the Urdu for a pole, then say \emph{dikhānā} again
\end{itemize}

\begin{tcolorbox}[breakable,colback=teal!4,colframe=teal!35!black,title={Before you move on}]
What does \ur{دکھانا} mean? (\textquotedblleft{}To show\textquotedblright{}.) Say it once more.
\end{tcolorbox}
